%%=============================================================================
%% Inleiding
%%=============================================================================

\chapter{\IfLanguageName{dutch}{Inleiding}{Introduction}}%
\label{ch:inleiding}

% De inleiding moet de lezer net genoeg informatie verschaffen om het onderwerp
% te begrijpen en in te zien waarom de onderzoeksvraag de moeite waard is om te
% onderzoeken. In de inleiding ga je literatuurverwijzingen beperken, zodat de
% tekst vlot leesbaar blijft. Je kan de inleiding verder onderverdelen in
% secties als dit de tekst verduidelijkt. Zaken die aan bod kunnen komen in de
% inleiding~\autocite{Pollefliet2011}:
%
% \begin{itemize}
%   \item context, achtergrond
%   \item afbakenen van het onderwerp
%   \item verantwoording van het onderwerp, methodologie
%   \item probleemstelling
%   \item onderzoeksdoelstelling
%   \item onderzoeksvraag
%   \item \ldots
% \end{itemize}

Forcebit Ticketing vertrekt vanuit een concreet probleem bij Forcebit: supportcommunicatie met klanten stond vooral verspreid over telefoon en e-mail. Daardoor was het niet altijd duidelijk welke vragen nieuw waren, welke al behandeld werden en welke bijlagen bij welke vraag hoorden. De applicatie centraliseert die communicatie in tickets met een titel, categorie, onderwerp, status, conversatie en bijlagen.

\section{\IfLanguageName{dutch}{Probleemstelling}{Problem Statement}}%
\label{sec:probleemstelling}

% Uit je probleemstelling moet duidelijk zijn dat je onderzoek een meerwaarde
% heeft voor een concrete doelgroep. De doelgroep moet goed gedefinieerd en
% afgelijnd zijn.

Een ticketingsysteem brengt structuur in supportcommunicatie, maar alleen wanneer de regels duidelijk blijven. Rechten, antwoorden, statusovergangen en bijlagen mogen niet verspreid raken over controllers, schermen of losse hulpfuncties. De uitdaging is dus niet alleen een werkende applicatie bouwen, maar ook een codebase die onderhoudbaar blijft wanneer de functionaliteit groeit.

\section{\IfLanguageName{dutch}{Onderzoeksvraag en doelstelling}{Research question and objective}}%
\label{sec:onderzoeksvraag}
\label{sec:onderzoeksdoelstelling}

% Wees zo concreet mogelijk bij het formuleren van je onderzoeksvraag. Een
% onderzoeksvraag is trouwens iets waar nog niemand op dit moment een antwoord
% heeft (voor zover je kan nagaan). Je kan de onderzoeksvraag verder
% specificeren in deelvragen.

\begin{quote}\small
Hoe kan een onderhoudbare ticketingapplicatie gebouwd worden waarin klant en administrator een duidelijke supportworkflow volgen, met aandacht voor Single Responsibility, gelaagde structuur en herbruikbare frontendlogica?
\end{quote}

Het doel is een werkend proof of concept te bouwen waarin:

\begin{itemize}[nosep, leftmargin=*]
  \item klanten tickets kunnen aanmaken, beantwoorden, sluiten en heropenen;
  \item administrators tickets kunnen zoeken, filteren, beantwoorden en van status veranderen;
  \item backendlagen, DTO's, domeinregels, repositories, hooks en herbruikbare componenten de onderhoudbaarheid ondersteunen.
\end{itemize}

Daarna bespreekt het verslag achtereenvolgens theorie, aanpak, functionele analyse, backend, databank, frontend, aanvullende functionaliteiten en conclusie.
